% Section 7: Discussion
% \input-able LaTeX file for Paper 5
% Conventions: a \sp b = sequential product (dot notation matching vdW)
%   C_p = Alfsen-Shultz compression (P-projection)
%   Axiom source: arXiv:1803.11139 Definition 2 EXCLUSIVELY

\section{Discussion}\label{sec:discussion}

%--------------------------------------------------------------------
\subsection{Summary of Results}\label{subsec:summary}

Self-modeling systems, as formalized in
Definition~\ref{def:self-modeling-system}, are exactly the complex
matrix algebras $\Mnsa{n}$ ($n \geq 2$) equipped with the L\"uders
sequential product
$\seqp{a}{b} = \sqrt{a}\, b\, \sqrt{a}$ and the conjugate-transpose
involution $X^{*} = X^{\dagger}$.  The identification proceeds
through a chain of eight steps (see \Cref{fig:chain}): the
self-modeling definition entails a sequential product on the
effect space; the product satisfies the seven axioms of van de
Wetering, establishing Euclidean Jordan algebra structure; only
complex systems admit the minimal composite the definition
requires; and the surviving algebra is a $C^{*}$-algebra via
published theorems of van de
Wetering~\cite{vandeWetering2019b},
Barnum--Wilce~\cite{BarnumWilce2014}, and
Hanche-Olsen~\cite{HancheOlsen1985}.

The complex field, the Jordan product, and the $C^{*}$-involution are
consequences of Definition~\ref{def:self-modeling-system}.  Local
tomography follows from Condition~\ref{sms:minimal} (minimal
composite) via state separation (lower bound) and minimality (upper
bound); the complex field and $C^{*}$-structure then follow from the
classification machinery.  None are independent postulates.
The four conditions of the definition each unpack one word of
``self-modeling system''; we now examine their individual roles,
restrictiveness, and the consequences of weakening them.

\paragraph{Dependency audit.}
We itemize explicitly what each ingredient contributes to the
conclusion, to prevent blurring of the derivation structure.
\begin{itemize}[nosep]
\item \textbf{Ambient framework}
  (spectral OUS with Alfsen--Shultz compressions and Peirce
  decomposition): provides the setting in which both classical and
  quantum theories live; supplies spectral decomposition, faces,
  compressions, and the Peirce direct sum.  This is rich structure
  but is \emph{not biased toward quantum mechanics}
  (Remark~\ref{rem:proximity-objection}).
\item \textbf{Definition~\ref{def:self-modeling-system}(i)--(ii)}
  (spectral OUS + faithful tracking):
  provides the operational product $\seqp{p}{q} = \comp{p}(q)$ on
  sharp effects and the Peirce $1$-space
  ansatz~\eqref{eq:general-product}; constrains the mixing function
  to $f = \sqrt{\lambda_i\, \lambda_j}$ via S5
  (Proposition~\ref{prop:coherence}).
  This is the novel construction of the paper.
\item \textbf{Published theorem}
  (van de Wetering~\cite{vandeWetering2019b}, Theorem~1):
  converts S1--S7 into EJA structure.
\item \textbf{Definition~\ref{def:self-modeling-system}(iii)--(iv)}
  (minimal composite + simplicity):
  yields local tomography (Theorem~\ref{thm:local-tomo}), which
  excludes all non-complex simple EJA types.
\item \textbf{Published theorems}
  (van de Wetering~\cite{vandeWetering2019b} Theorem~3,
  Barnum--Wilce~\cite{BarnumWilce2014} Theorem~4.1,
  BGW~\cite{BarnumUdudeckVdW2020} Theorem~1):
  promote the surviving EJA to $C^{*}$-algebra structure and
  exclude spin factors and the Albert algebra.
\end{itemize}
\noindent
The novel mathematics of the paper is concentrated in the second
item: the self-modeling sequential product construction and the proof
that S5 forces scalar Peirce $1$-space action
(Proposition~\ref{prop:coherence}).  Everything upstream is framework;
everything downstream is published classification machinery.


%--------------------------------------------------------------------
\subsection{The Four Conditions}\label{subsec:assumptions}

We analyze each condition's role, restrictiveness, and the
consequences of weakening it.

\paragraph{Condition~\ref{sms:finite}: Finite-dimensional spectral OUS.}
We argue in \Cref{subsec:unified} that this follows from finite
capacity of the embedding region.  It is the most restrictive
assumption in the formal sense.  The spectral decomposition
$a = \sum_{i} \lambda_{i}\, p_{i}$, which underlies the entire
product construction~\eqref{eq:general-product}, requires a
well-behaved spectral theory.  In finite dimensions, this is
guaranteed by the Alfsen--Shultz theory~\cite{AlfsenShultz2003}.  In
infinite dimensions, spectral theory introduces domain issues: the
spectral measure replaces the finite sum, unbounded effects arise, and
the Peirce decomposition becomes a conditional-expectation-valued
integral.  We expect the core argument to extend to type~I factors
(direct integrals of matrix algebras), where the spectral
decomposition retains its essential features.  Extension to type~II
or~III von~Neumann algebras would require substantially new techniques,
particularly for the Peirce feedback construction and the facial
orthogonality argument for~\ref{ax:S4}.

\paragraph{Condition~\ref{sms:faithful}: Faithful self-model.}
Faithfulness means $\varphi: V_{B} \to V_{M}$ is an order
isomorphism---the model is a perfect internal copy.  As discussed in
Remark~\ref{rem:division}, the mixing function
$f = \sqrt{\lambda_{i}\, \lambda_{j}}$ is determined by
internal consistency (\ref{ax:S5}), not directly by~$\varphi$.
Heuristically, $f = 0$ corresponds to a trivial (non-faithful) model
and $f = \sqrt{\lambda_{i}\, \lambda_{j}}$ to a faithful one, but we
do not prove a quantitative theorem linking the fidelity of~$\varphi$
to the value of~$f$.  Whether approximate self-models (where
$\varphi$ is injective but not surjective) yield ``partially
coherent'' products with $|f| < \sqrt{\lambda_{i}\, \lambda_{j}}$
is a conjectural extension that we leave to future work.

\paragraph{Condition~\ref{sms:minimal}: Minimal composite.}
We argue in \Cref{subsec:unified} that minimality is the natural
composite for a maximally efficient self-modeler.
The body--model composite $V_{BM}$ is defined as the minimal order
unit space carrying product states, product effects, non-signaling
constraints, and a product-form sequential product.  For complex
quantum mechanics, the minimal and maximal composites coincide: the
Hilbert space tensor product $M_{nm}(\mathbb{C})^{\mathrm{sa}}$ is
the unique composite satisfying
(C1)--(C4)~\cite{Barnum2023}.  This is a
feature of the complex case, not a coincidence: it means Condition~\ref{sms:minimal}
is not restrictive for the theory we derive.
For a priori exploration, one could ask whether maximal composites
yield different conclusions.  For the real and quaternionic cases,
minimal $\neq$ maximal: the difference manifests as a failure of local
tomography (the real composite has ``hidden entangled states'' that
product measurements cannot resolve).  For the non-complex types, the discrepancy between minimal and
maximal composites is exactly the opacity that
Section~\ref{subsec:unified} rules out on self-modeling grounds.

\paragraph{Condition~\ref{sms:simple}: Simplicity.}
We argue in \Cref{subsec:unified} that simplicity corresponds to
``whole models whole''---the self-modeler models all of itself, not
independent parts.
The trace form non-degeneracy argument in
\Cref{prop:nondeg} requires the EJA to be
simple---having no nontrivial Jordan ideal.  An EJA with a direct sum
decomposition $V = V_{1} \oplus V_{2}$ splits into superselection
sectors invisible to each other; a system in one sector cannot model
another, so the self-model does not cover the whole.  However, each simple
summand satisfies our theorem independently: if $V_{1}$ and $V_{2}$
each admit faithful self-models, they are each $\Mnsa{n_{i}}$ by the
main theorem applied to each summand.  The direct sum then describes a
system with superselection sectors, each governed by complex quantum
mechanics.  Thus Condition~\ref{sms:simple} restricts to irreducible
systems, which is the generic case.


%--------------------------------------------------------------------
\subsection{The Definition Revisited}\label{subsec:unified}

Definition~\ref{def:self-modeling-system} formalizes self-modeling as
four conditions.  Section~\ref{sec:self-modeling} argued that each
unpacks one word of ``self-modeling system.''  We now examine why each
clause is forced---and why removing any one would break the concept.

\paragraph{Finite capacity.}
A self-modeler that contains an isomorphic copy of itself within itself
must have finite state-space dimension.  An infinite-dimensional system
would require an infinite-dimensional internal model, but a subsystem
embedded in a finite spatial region has finite information
capacity~\cite{Bekenstein1981}.  The spectral structure
(distinguishable states, resolutions of unity) is not an abstract
regularity condition: it is the minimal apparatus for
\emph{probing}---a system that tests its own state requires states to
distinguish.  In this sense, Condition~\ref{sms:finite} is not a
restriction on the formalism but a structural consequence of what ``can
probe itself in a finite region'' means operationally.

\paragraph{Faithfulness.}
This is the operational premise itself.  No further argument is needed.

\paragraph{Internality.}
Condition~\ref{sms:minimal} (minimal composite) says the body--model
composite carries no structure beyond what product measurements can
resolve.  The argument is not efficiency but \emph{accessibility}.

A self-modeler's only interface with its own body--model composite is
through product measurements: test an effect on the body, test an
effect on the model, combine the outcomes.  If the composite carried
states that product measurements could not distinguish---an entangled
sector invisible to product observations---those states would be
aspects of the system's own composite that it cannot probe.  A system
with such opaque internal structure is not modeling \emph{itself}; it
is modeling a shadow of itself, blind to correlations it cannot
access.

In the non-complex EJA types, this opacity is exactly what occurs.
For real quantum mechanics ($n \ge 2$), the maximal composite has
$\dim > d^{2}$: there are body--model states that product
measurements cannot resolve, giving the self-model blind
spots~\cite{BarnumWilce2014}.  For quaternionic systems, $\dim <
d^{2}$: product effects are linearly dependent, making the
self-model's observations redundant.  In both cases, the self-modeler
fails to faithfully track its own composite state.  Only in the
complex case do the minimal and maximal composites
coincide~\cite{Barnum2023}, meaning product measurements see
\emph{everything}---the composite has no hidden structure, and the
self-model is genuinely faithful at the composite level.

Minimality is therefore not an independent design choice imposed from
outside.  It is forced by the conjunction of faithfulness
(\ref{sms:faithful}) and internality: if the only available
measurements are product measurements, and the model must faithfully
capture all accessible structure, then the composite can contain
nothing that product measurements cannot resolve.

\paragraph{Irreducibility.}
Condition~\ref{sms:simple} (simple EJA) says the algebra has no
nontrivial ideal---it cannot be decomposed into independent blocks.
A self-modeler that decomposes into independent subsystems is not
modeling \emph{all} of itself: each block models only itself, not the
whole.  Simplicity is the algebraic expression of ``whole models
whole.''  As noted in \Cref{subsec:assumptions}, each simple summand
independently satisfies the main theorem, so the assumption restricts
to irreducible systems rather than excluding composite ones.

\medskip\noindent
The mathematical theorem requires the four conditions as stated.
The conceptual argument is that they are not independent postulates
but the structural shape that self-modeling takes in a finite-capacity
system---the minimum a self-modeler could be.  The following
characterization makes the relationship precise.

\begin{corollary}[Characterization]\label{cor:equivalence}
  Among simple finite-dimensional EJAs, a type admits a faithful
  self-model whose minimal composite is locally tomographic if and
  only if it is $\Mnsa{n}$ for some $n \geq 2$.
\end{corollary}

\begin{proof}
  The forward direction is the main theorem.  For the converse,
  take $V_{M} = V_{B} = \Mnsa{n}$ with $\varphi = \mathrm{id}$.
  The minimal composite of $\Mnsa{n}$ with itself is
  $M_{n^{2}}(\mathbb{C})^{\mathrm{sa}}$, which satisfies
  $\dim = n^{4} = (n^{2})^{2} = [\dim(V_{B})]^{2}$.  Local
  tomography holds, and the product-form sequential product inherits
  S1--S7 from the factors.
\end{proof}

\begin{remark}[Why the other types fail]
  The converse fails for every non-complex simple EJA type.
  For real matrices ($n \ge 2$), the minimal composite has
  $\dim > [\dim(V)]^{2}$: product measurements leave an entangled
  sector unresolved, so the self-model has \emph{blind spots} about
  body--model correlations.  For quaternionic matrices ($n \ge 2$),
  $\dim < [\dim(V)]^{2}$: product effects are linearly dependent,
  so the self-model's observations are \emph{redundant}.  For spin
  factors and the Albert algebra, no locally tomographic composite
  exists at all.  Complex quantum mechanics is the unique theory in
  which a finite-dimensional system can have complete self-knowledge
  through product measurements on the body--model composite.
\end{remark}

\noindent
Corollary~\ref{cor:equivalence} says that the conditions of our
theorem are not merely sufficient for complex quantum mechanics---they
\emph{characterize} it.  There is no simple finite-dimensional EJA
satisfying our conditions that is not complex QM, and no complex QM
system that fails to satisfy them.  The ``axiomatic corridor''
identified by the theorem is, in this precise sense, the corridor of
self-modeling itself.


%--------------------------------------------------------------------
\subsection{Comparison with Other Programs}\label{subsec:comparison}

\Cref{tab:comparison} provides a summary comparison.  We now discuss
several programs in more depth.

\paragraph{Masanes--M\"uller (2011--2013).}
The Masanes--M\"uller program~\cite{MasanesMueller2011,MasanesMueller2013}
achieves the closest rival in parsimony, using three postulates:
the existence of an information unit, continuous reversibility, and
tomographic locality.  The key difference is that their tomographic
locality postulate assumes local tomography (effectively selecting the
complex field as input), whereas we \emph{derive} local tomography
from faithful tracking via the EJA trace form non-degeneracy argument
(\Cref{thm:local-tomo}).  Their continuous reversibility axiom---that
every reversible transformation can be continuously connected to the
identity---has no analogue in our framework; instead, the
reversibility structure emerges from the $C^{*}$-algebra promotion.
On the other hand, their framework handles infinite dimensions more
naturally than ours, and their postulates have a cleaner operational
interpretation in terms of information processing.

\paragraph{Chiribella--D'Ariano--Perinotti (2011).}
The CDP program~\cite{CDP2011} uses six axioms, with Purification as
the central workhorse: every mixed state of a system arises as the
marginal of a pure state of a larger system.  Purification is a
powerful constraint that immediately selects complex quantum mechanics
from the landscape of GPTs.  We derive the $C^{*}$-algebraic structure
without assuming purification; purification then follows as a
\emph{theorem} of $\Mnsa{n}$, not a postulate.  However, the CDP
framework provides a complete categorical description of quantum
theory (including channels, instruments, and higher-order maps) that
goes well beyond our algebraic result.

\paragraph{Van de Wetering's effectus theory (2019).}
Our work relies heavily on van de Wetering's
theorems~\cite{vandeWetering2019b}, but the relationship to his
effectus-theoretic reconstruction~\cite{vandeWetering2019} is
complementary rather than subsumptive.  The effectus approach uses
approximately five axioms on the categorical structure of a theory
(effect algebra homomorphisms, images, compressions, quotients, sharp
effects).  Our approach is more operational: we replace categorical
axioms with the self-modeling definition and derive the
sequential product structure that feeds into van de Wetering's
classification theorems.  The two programs share a common downstream
engine (Theorems~1 and~3 of~\cite{vandeWetering2019b}) but differ in
what they take as input.

\paragraph{Alfsen--Shultz geometry.}
The Alfsen--Shultz program~\cite{AlfsenShultz2003} characterizes
$C^{*}$-algebras via the geometry of state spaces, using the concept
of an \emph{orientation} on the state space.  Their characterization
is geometrically elegant: a JB-algebra is the self-adjoint part of a
$C^{*}$-algebra if and only if its state space admits a consistent
orientation (a choice of ``which way around'' each face).  In our
framework, the orientation is not an input but emerges from the
self-modeling construction: the tracking map $\varphi$ provides a
canonical correspondence between faces of $V_{B}$ and faces of $V_{M}$
that implicitly defines an orientation.  Making this connection precise
is a direction for future investigation.


%--------------------------------------------------------------------
\subsection{Future Directions}\label{subsec:future}

\paragraph{Infinite-dimensional extension.}
The most important open problem is extending the result to
infinite-dimensional systems.  The key obstacles are:
\begin{enumerate}[label=(\roman*)]
  \item the spectral decomposition becomes a spectral measure, and the
    Peirce feedback
    formula~\eqref{eq:general-product} becomes an operator-valued
    integral;
  \item the facial orthogonality argument for \ref{ax:S4} relies on
    Alfsen--Shultz propositions that may require modification in
    infinite dimensions;
  \item the minimal composite construction must be replaced by an
    appropriate operator-algebraic tensor product.
\end{enumerate}
We conjecture that the result extends to type~I factors, and that the
self-modeling condition for type~III algebras is vacuous (no faithful
finite-dimensional self-model exists).

\paragraph{Approximate self-models and decoherence.}
If the tracking map $\varphi$ is only approximately faithful (e.g.,
$\|\varphi(a) - a\| \leq \varepsilon\|a\|$ for small
$\varepsilon$), the mixing function $f$ would not saturate the
positivity bound.  The resulting ``partially coherent'' sequential
product interpolates between classical ($f = 0$, full decoherence) and
quantum ($f = \sqrt{\lambda_{i}\lambda_{j}}$, full coherence).  This
suggests a connection between self-modeling fidelity and decoherence:
a system that imperfectly models itself would exhibit intrinsic
decoherence proportional to the modeling error.  Formalizing this
connection could provide an operational account of the
quantum-to-classical transition.

\paragraph{Categorical formulation.}
The self-modeling construction could be formulated in the language of
effectus theory or process
theories~\cite{SelbyScandaloCoecke2021}, where the tracking
map $\varphi$ becomes a morphism in a suitable category and the
product construction becomes a functor.  This would clarify the
relationship to the categorical reconstruction programs and might
reveal additional structure (e.g., whether self-modeling systems form a
subcategory of effectus algebras with special properties).

\paragraph{Connection to stochastic--quantum bijection.}
Barandes~\cite{Barandes2025} has established a bijection between
stochastic processes and quantum dynamics, showing that unitary quantum
evolution can be recast as a stochastic process with indivisibility
constraints.  If a faithful self-model generates a stochastic process
(via the sequential test--update--test cycle), the Barandes bijection
would map it to a quantum channel.  Investigating whether the
self-modeling premise can be rephrased in stochastic terms---and
whether the Barandes bijection provides an alternative derivation
route---is an intriguing direction.

\paragraph{Non-Markovian self-models.}
Our construction assumes a memoryless update cycle: testing $p$
triggers a single compression $\comp{\varphi(p)}$, with no dependence
on prior measurements.  A non-Markovian self-model, where the update
depends on a history of prior tests, would introduce memory effects
and might connect to quantum memory channels or process matrices.


%--------------------------------------------------------------------
\subsection{Anticipated Objections}\label{subsec:objections}

We address several objections that a careful reader might raise.

\begin{remark}[``Isn't $\sqrt{\cdot}$ importing Hilbert space structure?'']
\label{rem:sqrt-objection}
No.  The square root in the mixing function operates on \emph{real
eigenvalues} $\lambda_{i}, \lambda_{j} \in [0,1]$: it is the
elementary real-valued function $\sqrt{x} : [0,1] \to [0,1]$, the
unique non-negative real number whose square is $x$.  No operator
square root, no spectral theorem for self-adjoint operators on a
Hilbert space, no functional calculus of $C^{*}$-algebras, and no
trace operation is invoked.  The distinction is between a scalar
function on the reals (which we use) and an operator function on a
$C^{*}$-algebra (which we do not).  See \Cref{sec:circularity} for a
complete inventory of the mathematical objects used in the
construction.
\end{remark}

\begin{remark}[``Isn't minimality just local tomography in disguise?'']
\label{rem:minimality-objection}
The objection is that Condition~\ref{sms:minimal} (minimal
composite) does the real work of excluding non-complex types, and
that calling it ``self-modeling'' is relabeling.  We take this
objection seriously and answer it on three levels.

\emph{Conceptual.}
The argument of Section~\ref{subsec:unified} (Internality) is that a
self-modeler whose only interface with its own composite is product
measurements \emph{cannot} have a composite carrying states those
measurements fail to resolve.  Minimality follows from faithfulness
(\ref{sms:faithful}) plus internality: the composite must contain
nothing the self-modeler cannot probe.  This is a constraint
\emph{about self-modeling}, not a free-standing information-processing
postulate.

\emph{Mathematical.}
For the theory we derive (complex quantum mechanics), the minimal
and maximal composites coincide---this is a theorem, not an
assumption~\cite{Barnum2023}.  So Condition~\ref{sms:minimal} is
\emph{automatically satisfied} by every complex system: it imposes
no restriction beyond what complex QM already has.  It restricts
only the non-complex types, where the discrepancy between minimal
and maximal composites manifests as exactly the blind spots or
redundancies that Section~\ref{subsec:unified} rules out on
self-modeling grounds.

\emph{Structural.}
Local tomography is a property of a composite
($\dim V_{BM} = (\dim V)^2$).  Minimality is a construction
principle for what the composite \emph{is}.  The theorem that the
minimal composite of a simple EJA with itself is locally tomographic
if and only if the EJA is complex is a \emph{consequence} of the
definition, not a restatement.  The derivation---state separation
for the lower bound, explicit construction of~$W$ for the upper
bound---is given in full in Theorem~\ref{thm:local-tomo}.
\end{remark}

\begin{remark}[``Isn't $\varphi$ mathematically inert?'']
\label{rem:phi-inert-objection}
The compression product $\seqp{p}{q} = \comp{p}(q)$ for sharp
effects exists in any spectral OUS with compressions, independent
of~$\varphi$.  However, compressions alone give only a partial
product (defined on sharp effects); the extension to general effects
requires the Peirce $1$-space
restoration~\eqref{eq:general-product}, whose operational
justification is the self-modeling feedback loop (the model updates,
then the body is re-tested).  Without the self-modeling cycle, one
has compressions but no reason to assemble them into a sequential
product on all effects.  The tracking map~$\varphi$ provides the
product; the axioms pin its parameters.  See
Remark~\ref{rem:division} for the division of labor.
\end{remark}

\begin{remark}[``What about infinite dimensions?'']
\label{rem:infty-objection}
Our result is restricted to finite-dimensional systems
(\Cref{ass:finite}).  This is stated explicitly in the abstract, in
the main theorem (\Cref{thm:main}), and in the discussion above.  The
finite-dimensional restriction is genuine, not cosmetic: the spectral
decomposition, the Peirce feedback construction, and the facial
orthogonality argument all rely on finite-dimensional spectral theory.
We view the infinite-dimensional extension as the most important open
problem; see \Cref{subsec:future}.
\end{remark}

\begin{remark}[``How does this compare to effectus theory?'']
\label{rem:effectus-objection}
Van de Wetering's effectus-theoretic
reconstruction~\cite{vandeWetering2019} uses categorical axioms on
effect algebras to arrive at the same classification theorems we
invoke.  The relationship is complementary: we replace categorical axioms
with the self-modeling definition
(Definition~\ref{def:self-modeling-system}) and derive the
sequential product structure that effectus theory axiomatizes
directly.  Our approach trades categorical generality for
operational specificity.
Both programs ultimately rely on the same downstream theorems
(Theorems~1 and~3 of~\cite{vandeWetering2019b}), so they are
complementary rather than competing.
\end{remark}

\begin{remark}[``Isn't the spectral OUS framework already 80\% quantum?'']
\label{rem:proximity-objection}
A spectral order unit space with compressions, Peirce decomposition,
and resolution of unity is rich structure---but it is not biased toward
quantum mechanics.  Classical probability theory lives comfortably
within this framework as the special case where all Peirce $1$-spaces
are trivial ($V_1(p_i, p_j) = \{0\}$ for all projective units) and
the sequential product reduces to pointwise multiplication on a
simplex.  The content of our derivation is that self-modeling \emph{requires}
non-trivial Peirce $1$-space content: the off-diagonal structure that
distinguishes quantum from classical.  Without self-modeling, one
could have a perfectly valid spectral OUS that is entirely classical.
The Alfsen--Shultz framework is the minimal setting in which both
classical and quantum theories can be formulated on equal footing;
self-modeling \emph{is} the quantum case.
\end{remark}

\begin{remark}[``Is faithful self-modeling physically motivated?'']
\label{rem:motivation-objection}
Self-modeling is motivated by the observation that quantum systems are
routinely used to model other quantum systems (quantum simulation, quantum
error correction, quantum tomography).  The faithfulness condition---that
the model captures all operationally accessible structure---is the
requirement that the simulation is \emph{exact}, not approximate.  In
quantum computation, a universal quantum computer can faithfully
simulate any quantum system of equal or smaller dimension; our premise
asks only that the system can simulate \emph{itself}.  Our result
is a mathematical theorem about order unit spaces, independent of
interpretive commitments about consciousness or cognition.
\end{remark}


%--------------------------------------------------------------------
\subsection{Conclusion}\label{subsec:conclusion}

Faithful self-modeling is complex quantum mechanics.  The four
conditions of Definition~\ref{def:self-modeling-system}---spectral
state space, faithful tracking, minimal composite, and
irreducibility---characterize exactly $\Mnsa{n}$.  Starting from
these conditions, we constructed a sequential product on the effect
algebra.  This product satisfies van de Wetering's seven axioms,
establishing Euclidean Jordan algebra structure.  The self-modeling
definition implies local tomography for the body--model composite,
which excludes all non-complex EJA types and, via the theorems of
van de Wetering, Barnum--Wilce, and Hanche-Olsen, promotes the algebra
to $\Mnsa{n}$ with the conjugate-transpose involution.

The complex field, the Jordan product, the $C^{*}$-involution, and
local tomography are not independent postulates of quantum theory.
They are consequences of the self-modeling definition
(Definition~\ref{def:self-modeling-system}).  The four conditions of
the definition---finite dimensionality, faithfulness, minimal
composite, and simplicity---unpack four words of ``self-modeling system'' and constitute what
self-modeling \emph{is} in the finite-dimensional OUS setting.

\paragraph{Machine verification.}
The derivation chain is formalized in Lean~4 with Mathlib (v4.28.0).
The formalization comprises 11 files (${\sim}3{,}400$ lines) and
produces zero \texttt{sorry} declarations.  Sixteen axioms are used,
all citing published external results: 14 encode order unit space
and Peirce decomposition properties from Alfsen--Shultz (2003,
Chapters~6--9), one encodes the spectral Jordan identity (van de
Wetering 2019, \S4), and one encodes the $C^{*}$-promotion theorem
(van de Wetering 2019, Theorem~3).  The novel construction of
Sections~3--4 (self-modeling to sequential product, all 10 SPData
fields) is \emph{fully proved}, not axiomatized.
Two concrete models---diagonal $2 \times 2$ matrices and the spin
factor---satisfy all S1--S7 from scratch, verifying that the axiom
system is consistent.  The formalization is available at
\url{https://github.com/ehrlich-b/research/tree/jmp-submission-2026-03-28/lean}.
